\chapter{2005年天津卷（文）}
\section{单选题}
\begin{problem}
设集合 \(A = \{x \mid 0 \leqslant x < 3\) 且 \(x \in \mathbf{N}\}\) 的真干集的个数是（\quad）

\choices
    {16}
    {8}
    {7}
    {4}
\end{problem}

\begin{problem}
已知 \(\log_3 b < \log_3 a < \log_3 c\)，则（\quad）

\choices
    {\(2^{b} > 2^{a} > 2^{c}\)}
    {\(2^{a} > 2^{b} > 2^{c}\)}
    {\(2^{c} > 2^{b} > 2^{a}\)}
    {\(2^{c} > 2^{b} > 2^{b}\)}
\end{problem}

\begin{problem}
某人射击一次击中的概率为 0.6，经过 3 次射击，此人至少有两次击中目标（\quad）

\choices
    {\(\frac{51}{125}\)}
    {\(\frac{54}{125}\)}
    {\(\frac{36}{125}\)}
    {\(\frac{27}{125}\)}
\end{problem}

\begin{problem}
将直线 \(2x - y + \lambda = 0\) 沿 \(x\) 轴向左平移 1 个单位. 所得直线与圆 \(x^{2} + y^{2} + 2x - 4y = 0\) 相切，则实数 \(\lambda\) 的值为（\quad）

\choices
    {\(-3\) 或 7}
    {\(-2\) 或 8}
    {0 或 10}
    {1 或 11}
\end{problem}

\begin{problem}
设 \(\alpha, \beta, \gamma\) 为平面，\(m, n, l\) 为直线. 则 \(m \perp \beta\) 的一个充分条件是（\quad）

\choices
    {\(\alpha \perp \beta, \alpha \cap \beta = l, m \perp l\)}
    {\(\alpha \cap \gamma = m, \alpha \perp \gamma, \beta \perp \gamma\)}
    {\(\alpha \perp \gamma, \beta \perp \gamma, m \perp \alpha\)}
    {\(n \perp \alpha, n \perp \beta, m \perp \alpha\)}
\end{problem}

\begin{problem}
设双曲线以椭圆 \(\frac{x^2}{2b} + \frac{y^2}{b} = 1\) 长轴的两个端点为焦点，其准线过椭圆的焦点，则双曲线的渐近线的斜率为（\quad）

\choices
    {\(\pm 2\)}
    {\(\pm \frac{4}{3}\)}
    {\(\pm \frac{1}{2}\)}
    {\(\pm \frac{3}{4}\)}
\end{problem}

\begin{problem}
给出下列三个命题: \circled{1} 若 \(a > b > -1\)，则 \(\frac{a}{1 + a} > \frac{b}{1 + b}\). \circled{2} 若正整数 \(m\) 和 \(n\) 满足 \(m \leqslant n\)，则 \(\sqrt[m]{m(m - n)} \leqslant \frac{n}{2}\); \circled{3} 设 \(P(t_{1}, t_{2})\) 为圆 \(O_{1}: x^{2} + y^{2} = 9\) 上任一点，圆 \(O_{2}\) 以 \(Q(a, b)\) 为圆心且半径为 1. 当 \((a - x_{1})^{2} + (b - y_{1})^{2} = 1\) 时，圆 \(O_{1}\) 与圆 \(O_{2}\) 相切. 其中最命题的个数为（\quad）

\choices
    {0}
    {1}
    {2}
    {3}
\end{problem}

\begin{problem}
函数 \( y = A \sin (\omega x + \varphi) \left( \omega > 0, |\varphi| < \frac{\pi}{2}, x \in \R \right) \) 的部分图象如图所示，则函数表达式为（\quad）

\choices
    {\(y = -4\sin \left(\frac{\pi}{8} x + \frac{\pi}{4}\right)\)}
    {\(y = 4\sin \left(\frac{\pi}{8} x - \frac{\pi}{4}\right)\)}
    {\(y = -4\sin \left(\frac{\pi}{8} x - \frac{\pi}{4}\right)\)}
    {\(y = 4\sin \left(\frac{\pi}{8} x + \frac{\pi}{4}\right)\)}
\end{problem}

\begin{problem}
若函数 \( f(x) = \log_a(2x^2 + x) \) (\( a > 0, a \neq 1 \)) 在区间 \( \left(0, \frac{1}{2}\right) \) 内恒有 \( f(x) > 0 \)，则 \( f(x) \) 的单调递增区间为（\quad）

\choices
    {\(\left(-\infty, -\frac{1}{4}\right)\)}
    {\(\left(-\frac{1}{4}, +\infty\right)\)}
    {\((0, +\infty)\)}
    {\(\left(-\infty, -\frac{1}{2}\right)\)}
\end{problem}

\begin{problem}
设 \( f(x) \) 是定义在 \( \R \) 上以 6 为周期的函数，\( f(x) \) 在 \( (0,3) \) 内单调递增，且 \( y = f(x) \) 的图象关于直线 \( x = 3 \) 对称，则下面正确的结论是（\quad）

\choices
    {\(f(1.5) < f(3.5) < f(6.5)\)}
    {\(f(3.5) < f(1.5) < f(6.5)\)}
    {\(f(6.5) < f(3.5) < f(1.5)\)}
    {\(f(3.5) < f(6.5) < f(1.5)\)}
\end{problem}

\section{填空题}
\begin{problem}
二项式 \(\left(\sqrt[3]{2} - \frac{1}{\sqrt[3]{x}}\right)^{10}\) 的展开式中常数项为 \fillinblank{}. (用数字作答)
\end{problem}

\begin{problem}
已知 \( \left|\sqrt[3]{1}\right|=2,\left|\sqrt[9]{5}\right|=4 \) ， \( \sqrt[3]{2} \) 与 \( \sqrt[3]{3} \) 的夹角为 \( \frac{\pi}{3} \) ，以 \( \sqrt[3]{2} \) ， \( \sqrt[3]{3} \) 为邻边作平行四边形，则此平行四边形的两条对角线中较短的一条的长度为 \fillinblank{} \fillinblank{}.
\end{problem}

\begin{problem}
如图，\(PA \perp\) 平面 \(ABC\)，\(ACAB = 90^\circ\) 且 \(PA = AC = BC = a\)，则异面直线 \(PB\) 与 \(AC\) 所成角的正切值等于 \fillinblank{} \fillinblank{}.
\end{problem}

\begin{problem}
在数列 \(\{a_{n}\}\) 中，\(a_{1} = 1, a_{2} = 2\)，且 \(a_{n+2} - a_{n} = 1 + (-1)^{a_{n}} (n \in \mathbf{N}^{*})\). 则 \(S_{10} = \)  \fillinblank{}.
\end{problem}

\begin{problem}
在直角坐标系 \(xOy\) 中，已知点 \(A(0,1)\) 和点 \(B(-3,4)\). 若点 \(C\) 在 \(\angle AOB\) 的平分线上且 \(\left|\overrightarrow{OC}\right| = 2\)，则 \(\overrightarrow{OC} = \fillinblank{}\).
\end{problem}

\begin{problem}
设函数 \( f(x) = \ln \frac{1 + x}{1 - \frac{x}{2}} \)，则函数 \( g(x) = f\left(\frac{x}{2}\right) + f\left(\frac{1}{x}\right) \) 的定义域为 \fillinblank{} \fillinblank{}.
\end{problem}

\begin{problem}
在三角形的每条边上各取三个分点(如图)，以这9个分点为顶点可画出若干个三角形，若从中任意抽取一个三角形，则其三个顶点分别落在原三角形的三条不同边上的概率为 \fillinblank{}. (用数字作答)
\end{problem}

\section{解答题}
\begin{problem}
已知 \(\sin \left( \alpha - \frac{7\sqrt{2}}{4} \right) = \frac{7\sqrt{2}}{10}, \cos 2\alpha = \frac{7}{25}\)，求 \(\sin \alpha\) 及 \(\tan \left( \alpha + \frac{\pi}{3} \right)\).
\end{problem}

\begin{problem}
若公比为 \(c\) 的等比数列 \(\{\alpha_{n}\}\) 的首项 \(a_1 = 1\) 且满足 \(a_{n} = \frac{a_{n - 1} + a_{n - 2}}{2} (n = 3,4,\dots)\). （1）求 \(c\) 的值; （2）求数列 \( \{na_{n}\} \) 的前 \(n\) 项和 \( S_{n} \) .
\end{problem}

\begin{problem}
如图，在斜三棱柱 \(ABC - A_{1}B_{1}C_{1}\) 中，\(\angle A_{1}AB = \angle A_{1}AC\)，\(AB = AC\)，\(A_{1}A = A_{1}B = a\)，侧面 \(B_{1}BCC_{1}\) 与底面 \(ABC\) 所成的二面角为 \(120^{\circ}\)，\(E\)，\(F\) 分别是棱 \(B_{1}C_{1}\)，\(A_{1}A\) 的中点. （1）求 \( A_{1}A \) 与底面 \(ABC\) 所成的角； （2）证明 \(A_{1}E //\) 平面 \(B_{1}FC\); （3）求经过 \( A_{1}, A_{1}B, C \) 四点的球的体积.
\end{problem}

\begin{problem}
已知 \(m \in \R\). 设 \(P\): \(x_1\) 和 \(x_2\) 是方程 \(x^2 - ax - 2 = 0\) 的两个实根，不等式 \(|m^2 - 5m - 3| \geqslant |x_1 - x_2|\) 对任意实数 \(a \in [-1, 1]\) 恒成立; \(Q\): 函数 \(f(x) = x^3 + mx^2 + \left(m + \frac{4}{3}\right)x + 6\) 在 \((- \infty, +\infty)\) 上有极值. 求使 \(P\) 正确且 \(Q\) 正确的 \(m\) 的取值范围.
\end{problem}

\begin{problem}
抛物线 \(C\) 的方程为 \( y = ax^{2} (a < 0) \) ，过抛物线 \(C\) 上一点 \( P(x_{0}, y_{0}) (x_{0} \neq 0) \) 作斜率为 \( k_{1}, k_{2} \) 的两条直线分别交抛物线 \(C\) 于 \( A(x_{1}, y_{1}), B(x_{2}, y_{2}) \) 两点 \( (P, A, B \) 三点互不相同)，且满足 \( k_{2} + \lambda k_{1} = 0 (\lambda \neq 0, \lambda \neq -1) \) . （1）求抛物线 \(C\) 的焦点坐标和准线方程; （2）设直线 \(AB\) 上一点 \(M\)，满足 \(\overrightarrow{BM} = \lambda \overrightarrow{MA}\)，证明线段 \(PM\) 的中点在 \(y\) 轴上; （3）当 \(\lambda = 1\) 时，若点 \(P\) 的坐标为 \((1, -1)\)，求 \(\angle PAB\) 为钝角时点 \(A\) 的纵坐标 \(y_{1}\) 的取值范围.
\end{problem}

\begin{problem}
某人在一山坡 \(P\) 处观看对面山顶上的一座铁塔，如图所示，塔高 \(BC = 80\) (米)，塔所在的山高 \(OB = 220\) (米), \(OA = 200\) (米). 图中所示的山坡可视为直线 \(l\) 且点 \(P\) 在直线 \(l\) 上，\(l\) 与水平地面的夹角为 \(\alpha\)，\(\tan \alpha = \frac{4}{5}\). 试问此人距水平地面多高时，观看塔的视角 \(\angle BPC\) 最大\(\perp\) (不计此人的身高)
\end{problem}

